% results_rq6

\section{RQ6: Completeness and the Two Gaps}
\label{sec:rq6}

The final research question asks whether current devotional apps are enough, in
the sense the ``apps under-serve niche needs'' framing implies. The answer is
no, but that word conceals two structurally different failures calling for
different remedies. A \emph{coverage gap} is a market failure: entire
categories of need are served by a minority of the ecosystem despite measurable
demand. A \emph{delivery gap} is a quality failure, and the more damning,
because it falls on promises apps have already made --- including the most
basic ones.

\subsection{RQ6a: The Delivery Gap}

We compared each application's store-description promises against the sentiment
its reviews record. Of 272 promised (application, feature) pairs, 136 carry
enough mentions to judge reliably (\(n \geq 30\)); of those, 18 are broken
promises across 11 applications --- a claimed feature drawing over 40\%
negative sentiment. Table~\ref{tab:rq6-broken} ranks the break rate by share of
promising applications and Figure~\ref{fig:rq6a-delivery-gap} shows the full
gap. The two largest by mention volume are both Muslim Pro's: reminders/adhan
(9,989 mentions, 45.6\% negative) and monetisation terms (9,182 mentions,
54.8\% negative).

\begin{table}[t]
  \centering
  \small
  \caption{Features with at least one broken promise among those with \(n \geq 30\) judgeable mentions and more than one promising application. \texttt{tracker\_score}, with a single promising application, is discussed below rather than listed.}
  \label{tab:rq6-broken}
  \begin{tabular}{lr}
    \toprule
    Feature & Broken / promising apps \\
    \midrule
    Monetisation             & 5 / 16 (31\%) \\
    Reminders / adhan        & 5 / 20 (25\%) \\
    Mosque finder             & 1 / 4 (25\%) \\
    Forbidden times            & 1 / 4 (25\%) \\
    Widgets                    & 2 / 13 (15\%) \\
    Prayer-time accuracy      & 2 / 19 (11\%) \\
    Qibla                      & 1 / 18 (6\%) \\
    \bottomrule
  \end{tabular}
\end{table}

Two rows carry a caveat: forbidden times is tagged at 10\% precision and mosque
finder at 27\% (\S\ref{sec:method-measurement}), so their 1-of-4 ratios rest on
few and noisy mentions, and we rest no argument on either.

\begin{figure}[t]
  \centering
  \includegraphics[width=\linewidth]{fig10_rq6a_delivery_gap}
  \caption{The delivery gap: negative-sentiment rate per (application, feature) pair across the 136 judgeable pairs.}
  \label{fig:rq6a-delivery-gap}
\end{figure}

One further aspect, \texttt{tracker\_score}, shows 1 of 1 promising
applications broken; with a single promising application the ratio carries no
statistical weight, so we do not present it as the most-broken feature.

The two features broken in the most applications, five each, are also the two
most basic. Adhan reminders are the core function of a prayer-times
application, and monetisation the terms on which it is offered. Failure here
concentrates in the promise the app is built around, not in peripheral
functionality.

\subsection{RQ6b: The Coverage Gap}

We separately ranked aspects by an unmet-need score --- demand-signal volume
multiplied by applications lacking the feature --- built from the
demand-extraction pipeline validated at 91.5\% precision
(\S\ref{sec:method-demand}) rather than from the noisier raw aspect counts
several low-precision aspects would otherwise force on us
(Table~\ref{tab:rq6-unmet}, Figure~\ref{fig:rq6b-unmet-needs}). Companion
hardware tops the ranking, followed by menstrual handling, mosque finder,
forbidden-times awareness, and \emph{nafl} (supererogatory) prayer-time
support. Four aspects fall below 25\% ecosystem coverage:
\texttt{calendar\_sync}, \texttt{companion\_hardware}, \texttt{qasr\_travel},
and \texttt{women\_period}.

\begin{table}[t]
  \centering
  \small
  \caption{Unmet-need ranking: demand-signal volume \(\times\) applications lacking the feature.}
  \label{tab:rq6-unmet}
  \begin{tabular}{lrr}
    \toprule
    Aspect & Score & Apps lacking \\
    \midrule
    Companion hardware & 6{,}725 & 25 \\
    Menstrual handling (\texttt{women\_period}) & 4{,}738 & 23 \\
    Mosque finder       & 3{,}304 & 14 \\
    Forbidden times      & 2{,}805 & 15 \\
    Nafl (supererogatory) prayer times & 2{,}445 & 15 \\
    \bottomrule
  \end{tabular}
\end{table}

\begin{figure}[t]
  \centering
  \includegraphics[width=\linewidth]{fig11_rq6b_unmet_needs}
  \caption{Unmet-need ranking across aspects: demand-signal volume against the number of applications lacking the feature.}
  \label{fig:rq6b-unmet-needs}
\end{figure}

\subsection{Some Needs Are Served Well}

Five aspects clear both bars we set for adequate service --- present in over
half the ecosystem, averaging under 30\% negative sentiment across the
applications where mentioned: madhab, monetisation, qibla, table format,
widgets. Table format is tagged at 44\% precision, so its place is the least
secure. The average is taken over applications rather than reviews, which
matters: under review-weighted pooling, monetisation, widgets, and qibla all
exceed 30\%, because a few high-volume applications handle them badly. We
report the app-level figure because the question is how many applications serve
a need adequately, not how many reviews are unhappy, but the choice should not
pass unnoticed.

Monetisation appears on both this list and the broken-promises list in
Table~\ref{tab:rq6-broken}, which is informative rather than contradictory.
Across the same 16 applications promising it with \(n \geq 30\) mentions, the
break rate is 5 of 16 (31\%) against a mean negative rate of 25.8\%: most
handle payment acceptably while a minority handle it badly enough to poison the
feature ecosystem-wide.

\subsection{Synthesis: The Same Gaps, Reached Independently}

Both gaps are also reached independently by an earlier research question, which
we read as corroboration rather than repetition. Companion hardware tops the
unmet-need ranking, and RQ5 finds hardware sentiment \emph{better} inside the
one application shipping a device than across the 23 that do not
(\S\ref{sec:rq5}) --- demand rather than dissatisfaction with an existing
product. Menstrual handling ranks second and is RQ4's central finding in its
own right (\S\ref{sec:rq4}). Reminders/adhan, tied with monetisation at five
broken applications each, is independently the dominant
tracker-dissatisfaction theme in RQ1's cross-app coding (\S\ref{sec:rq1}).

\subsection{Verdict}

Current applications are not enough, and fall short in two ways calling for
different responses. The coverage gap is a market failure: entire categories of
need --- companion hardware, menstrual accommodation, travel concession,
calendar integration --- are served by under a quarter of the ecosystem despite
measurable, explicitly voiced demand. The delivery gap is a quality failure and
the more damning, because it falls on promises already made: the features most
often broken are adhan reminders and monetisation, the two things a
prayer-times application most fundamentally offers. Building more features is
not indicated. Delivering the promised ones reliably, and extending coverage to
the populations currently unserved, is.
